\documentclass{article}
\usepackage{amsmath, amssymb}
\begin{document}

\section*{CSE 400: Fundamentals of Probability in Computing}

\textbf{Lecture 6: Discrete RVs, Expectation and Problem Solving}  
\textbf{Instructor:} Dhaval Patel, PhD  
\textbf{Date:} January 22, 2025  

\section*{1. Random Variables (RVs)}

\subsection*{1.1 Motivation and Concept}

A random variable $X$ on a sample space $\Omega$ is a function $X:\Omega\rightarrow\mathbb{R}$ that assigns a real number $X(\omega)$ to each sample point $\omega\in\Omega$.

\textbf{Discrete Random Variables (DRVs):} These take values in a range that is finite or countably infinite.

\textbf{Support:} Even though $X$ maps to $\mathbb{R}$, the actual set of values $\{X(\omega):\omega\in\Omega\}$ is a discrete subset of $\mathbb{R}$.

\subsection*{1.2 Example: Tossing 3 Fair Coins}

Let $Y$ denote the number of heads that appear. The possible values for $Y$ are $\{0, 1, 2, 3\}$.

\[
P\{Y=0\} = P\{(t,t,t)\} = \frac{1}{8}.
\]

\[
P\{Y=1\} = P\{(t,t,h),(t,h,t),(h,t,t)\} = \frac{3}{8}.
\]

\[
P\{Y=2\} = P\{(t,h,h),(h,t,h),(h,h,t)\} = \frac{3}{8}.
\]

\[
P\{Y=3\} = P\{(h,h,h)\} = \frac{1}{8}.
\]

\textbf{Total Probability:}
\[
\sum_{i=0}^{3}P\{Y=i\} = 1.
\]

\section*{2. Probability Mass Function (PMF)}

\subsection*{2.1 Definition}

Let $X$ be a discrete random variable with range
\[
R_{X}=\{x_{1},x_{2},x_{3},...\}.
\]
The function
\[
P_{X}(x_{k})=P(X=x_{k})
\]
for $k=1,2,3,...$ is called the Probability Mass Function (PMF) of $X$.

\subsection*{2.2 Properties of a Legitimate PMF}

\textbf{Probability Range:} $P(X=x_{k})$ gives the probability of each outcome.

\textbf{Normalization:}
\[
\sum_{k=1}^{\infty}P_{X}(x_{k})=1.
\]

\subsection*{2.3 Visualization}

The distribution can be visualized as a bar diagram where the x-axis represents possible values and the height of the bar at value $a$ is the probability $Pr[X=a]$.

\subsection*{2.4 Worked Example: Determining the Constant $c$}

Given a PMF
\[
p(i)=c\frac{\lambda^{i}}{i!}
\]
for $i=0,1,2,...$ where $\lambda > 0$.

\textbf{Summation Requirement:}
\[
\sum_{i=0}^{\infty}p(i)=1.
\]

\textbf{Derivation:}
\[
c\sum_{i=0}^{\infty}\frac{\lambda^{i}}{i!}=1.
\]

\textbf{Identity:} Using the Taylor series expansion
\[
e^{x}=\sum_{i=0}^{\infty}\frac{x^{i}}{i!},
\]
we substitute $x=\lambda$ to get
\[
ce^{\lambda}=1.
\]

\textbf{Result:}
\[
c=e^{-\lambda}.
\]

\textbf{Specific Probabilities:}
\[
P\{X=0\}=e^{-\lambda}\frac{\lambda^{0}}{0!}=e^{-\lambda}.
\]

\[
P\{X>2\}=1-P\{X\le2\}=1-P\{X=0\}-P\{X=1\}-P\{X=2\}.
\]

\section*{3. Bayes' Theorem and Independence}

\subsection*{3.1 Bayes' Formula Recap}

\[
Pr(B_{i}|A)=\frac{Pr(A|B_{i})Pr(B_{i})}{\sum_{j=1}^{n}Pr(A|B_{j})Pr(B_{j})}
\]

\textbf{A priori probability:} $Pr(B_{i})$ formed from presupposed models.

\textbf{Posteriori probability:} $Pr(B_{i}|A)$ derived after observing event $A$.

\subsection*{3.2 Independent Events}

Two events $A$ and $B$ are independent if
\[
Pr(A|B)=Pr(A) \quad \text{and} \quad Pr(B|A)=Pr(B).
\]

\textbf{Joint Probability:}
\[
Pr(A,B)=Pr(A)Pr(B).
\]

\textbf{Mutual Independence (3 Events):}
Events $A, B,$ and $C$ are mutually independent if each pair is independent and
\[
Pr(A,B,C)=Pr(A)Pr(B)Pr(C).
\]

\section*{4. Types of Discrete Random Variables}

\subsection*{4.1 Bernoulli Random Variable}

\textbf{Concept:} An experiment with two outcomes: Success (1) or Failure (0).

\textbf{PMF:}
\[
P_{X}(1)=p \quad \text{and} \quad P_{X}(0)=1-p.
\]

\textbf{Parameter:} $p \in (0,1)$ is the success probability.

\textbf{Examples:} Tossing a coin, email spam classification.

\subsection*{4.2 Binomial Random Variable $B(n,p)$}

\textbf{Concept:} The number of successes in $n$ independent trials, each with success probability $p$.

\textbf{PMF:}
\[
p(i)=\binom{n}{i}p^{i}(1-p)^{n-i}
\]
for $i=0,1,...,n$.

\textbf{Examples:} Correct answers in a test, number of defective items in a sample.

\subsection*{4.3 Geometric Random Variable}

\textbf{Concept:} Number of independent trials needed until the first success occurs.

\textbf{PMF:}
\[
P_{X}(n)=(1-p)^{n-1}p
\]
for $n=1,2,...$.

\textbf{Legitimacy Proof:}
\[
\sum_{n=1}^{\infty}p(1-p)^{n-1} = p \cdot \frac{1}{1-(1-p)} = 1.
\]

\textbf{Examples:} Tosses until the first head, attempts needed for a sale.

\subsection*{4.4 Poisson Random Variable}

\textbf{Concept:} Number of times a rare event occurs in a unit.

\textbf{PMF:}
\[
p(i)=e^{-\lambda}\frac{\lambda^{i}}{i!}
\]
for $i=0,1,2,...$ where $\lambda > 0$.

\textbf{Note:} Used as an approximation for Binomial $B(n,p)$ when $n$ is large and $p$ is small such that $np$ is moderate.

\textbf{Examples:} Misprints on a page, number of wrong telephone numbers dialed in a day.

\section*{5. Summary Table}

\begin{center}
\begin{tabular}{|l|l|l|}
\hline
Distribution & Description & Support \\
\hline
Bernoulli & Single trial (S/F) & $\{0,1\}$ \\
Binomial & \# of successes in $n$ independent trials & $\{0,1,\ldots,n\}$ \\
Poisson & \# of rare events in a unit & $\{0,1,2,\ldots\}$ \\
Geometric & \# of trials until 1st success & $\{1,2,3,\ldots\}$ \\
\hline
\end{tabular}
\end{center}

Would you like me to reconstruct the probability calculations for the ``Auditorium with 30 Rows'' example from the Bayes' Theorem section?

\end{document}
